\section{合卷：晋阳城下，春秋最后一个大国被分开}

晋曾经是春秋北方最有力量的国家。它能在城濮集结联盟，能在鄢陵与楚对抗，能让诸侯把会盟当作天下大事。可它的军队、封邑和政治资源从来不只在国君手里。卿族替国君打仗，也在战争、赏赐和封地中累积自己的家臣、土地和人脉。到春秋末年，晋国仍叫晋国，真正能决定晋国命运的人却越来越不是晋君。

\eventref{WQ-0453-ZHI}{战国·三家灭智}

前453年，智氏围赵于晋阳。后来的叙事把智伯写得骄横，把韩、赵、魏写成在危局中结盟的三家。围城的细节很有戏剧性：水、城、盟约、背叛都集中在一处。无论具体对白是否经过后人整理，有一件事无法回避：智氏试图用晋国旧有的卿族政治压倒赵氏，韩、赵、魏却发现，若各自等待，谁都逃不过被逐个吞并的命运。

于是三家合力灭智。智氏的土地、人口和军队被重新分割。这里发生的不是一个家族取代另一个家族的宫廷插曲，而是春秋大国的内部结构被拆开：原本服务于晋的资源，开始分别服务于韩、赵、魏。

\eventref{WQ-0403-THREEJIN}{战国·三家分晋}

五十年后，周威烈王册命韩、赵、魏为诸侯。王室仍然给出了名分，但名分来得很晚：三家早已拥有事实上的国家。周王室的册命像一纸迟到的文书，它没有创造新秩序，只是承认旧秩序已经无法恢复。

\begin{turningpoint}[三家分晋为什么是春秋和战国之间的门槛]
春秋时期，强国仍可以通过会盟、礼仪和霸主来维持一种可谈判的秩序；三家分晋之后，原本的大国被几支拥有独立军政资源的家族分割。国家竞争从“谁主持联盟”进一步变成“谁能直接掌握土地、户籍、军队和官吏”。魏、赵、韩会分别走向不同命运，但它们都从晋的遗产里带走了战国国家所需的原料。
\end{turningpoint}

\begin{sourcenote}[晋阳城下的故事有多少是史实]
关于智伯、晋阳围城和三家结盟的完整叙事，主要来自后出的《史记》《战国策》及《资治通鉴》所整理的材料。它们保留了强烈的政治戏剧性。这里不把每段劝说、每个水攻细节当作现场实录；叙事依赖的核心事实是，韩赵魏在前453年灭智并分其地，周王室于前403年承认三家诸侯地位。两件事之间的时间差，正是“事实上的国家”先出现、旧名分后追认的最好证据。
\end{sourcenote}
